\begin{table}[H]
\centering\small
\setlength{\tabcolsep}{4pt}
\renewcommand{\arraystretch}{1.12}
\caption{The 16 tools available to the agents and grouped into  seven families of task}\label{tab:retail-tools}
\begin{tabular}{@{}>{\raggedright\arraybackslash}p{.22\textwidth}>{\raggedright\arraybackslash}p{.75\textwidth}@{}}
\toprule Family suffix & Tool names\\\midrule

\feat{identity} & \feat{find_user_id_by_email}, \feat{find_user_id_by_name_zip}\\
\feat{user_read} & \feat{get_user_details}\\
\feat{order_read} & \feat{get_order_details}\\
\feat{catalog_read} & \feat{get_item_details}, \feat{get_product_details}, \feat{list_all_product_types}\\
\feat{pending_write} & \feat{cancel_pending_order}, \feat{modify_pending_order_address}, \feat{modify_pending_order_items}, \feat{modify_pending_order_payment}\\
\feat{resolution_write} & \feat{exchange_delivered_order_items}, \feat{modify_user_address}, \feat{return_delivered_order_items}\\
\feat{utility} & \feat{calculate}, \feat{transfer_to_human_agents}\\
\bottomrule
\end{tabular}
\end{table}
\begin{table}[H]
\centering\small
\setlength{\tabcolsep}{4pt}
\renewcommand{\arraystretch}{1.12}
\caption{The 24 Retail SNMM inputs, used for all continuations. Yes/no inputs
are one for yes and zero for no.}\label{tab:retail-snmm}
\begin{tabular}{@{}>{\raggedright\arraybackslash}p{.34\textwidth}r>{\raggedright\arraybackslash}p{.56\textwidth}@{}}
\toprule Input & Count & Meaning\\\midrule
\feat{intercept} & 1 & A constant value of one.\\
Proposed tools (\feat{tool:}) & 16 & One input for each tool in Table~\ref{tab:retail-tools}: whether the current proposal includes that tool.\\
\feat{prior_error} & 1 & An earlier tool call returned an error.\\
\feat{prior_correction} & 1 & An earlier review changed a proposal.\\
\feat{retry_tool} & 1 & At least one proposed tool was used earlier.\\
\feat{used:catalog_read} & 1 & An earlier product or catalog read succeeded.\\
\feat{used:order_read} & 1 & An earlier order-details read succeeded.\\
\feat{write_x_prior_error} & 1 & The proposal contains a write, and an earlier tool call returned an error.\\
\feat{write_x_affirmation} & 1 & The proposal contains a write, and an earlier user message contains a confirmation phrase.\\
\bottomrule
\end{tabular}
\end{table}

 A confirmation phrase matches
\texttt{yes|confirm|go ahead|proceed} as whole words, ignoring capitalization.
We count a tool result as an error if its error flag is set or its text starts
with \texttt{error}.

\begin{table}[H]
\centering\small
\setlength{\tabcolsep}{4pt}
\renewcommand{\arraystretch}{1.12}
\caption{The 41 features for the Retail nuisance models. }\label{tab:retail-nuisance}
\begin{tabular}{@{}>{\raggedright\arraybackslash}p{.34\textwidth}r>{\raggedright\arraybackslash}p{.56\textwidth}@{}}
\toprule Inputs & Count & Meaning\\\midrule
Proposed tools (\feat{tool:}) & 16 & The same 16 tool inputs as in Table~\ref{tab:retail-snmm}.\\
Current tool families (\feat{current:}) & 7 & For each family in Table~\ref{tab:retail-tools}, whether the proposal includes a tool from that family.\\
Earlier successful families (\feat{used:}) & 7 & For each family, whether an earlier call from that family succeeded.\\
Earlier events & 5 & Earlier error, earlier correction, and repeated tool (defined above); any earlier review assignment (\feat{prior_verification}); earlier user confirmation (\feat{user_affirmation_seen}).\\
\feat{n_prior_tools}, \feat{n_prior_errors}, \feat{n_prior_verifications}, \feat{n_proposed_tools} & 4 & Numbers of earlier tool results, earlier errors, earlier review assignments, and calls in the current proposal, respectively.\\
\feat{source_cohort_extension} & 1 & Whether the task comes from the added-task set.\\
\feat{eligible_position_zero_based} & 1 & The number of review opportunities before the current proposal.\\
\bottomrule
\end{tabular}
\end{table}

To predict terminal success, we also include the current review assignment
and earlier decision records. Each earlier record contains the 39 inputs in
the first five rows of Table~\ref{tab:retail-nuisance} and its review assignment.
We reserve 28 slots for GPT-5 nano and 14 for Claude Haiku; unused slots are zero.

The models used to account for later reviews use the 41 current inputs above,
plus one yes/no input for each possible decision number. Only the input for
the current decision number is one.
